\documentclass[11pt,a4paper]{article}
\usepackage[T1]{fontenc}
\usepackage[utf8]{inputenc}
\usepackage{amsmath,amssymb,xcolor}
\usepackage[margin=25mm]{geometry}
\usepackage[colorlinks=true,urlcolor=blue]{hyperref}
\definecolor{nexorablue}{HTML}{2459B5}
\setlength{\parindent}{0pt}
\setlength{\parskip}{7pt}
\setlength{\emergencystretch}{2em}
\renewcommand{\familydefault}{\sfdefault}
\begin{document}

{\small\bfseries\color{nexorablue} NEXORA CPGE}\par
{\small MPSI / NEX-RGT-05 / Énoncé}\par
{\LARGE\bfseries Régler une décharge au régime critique}\par
\vspace{0.25em}\hrule\vspace{0.6em}

Un circuit RLC série libre a $L=40\,\mathrm{mH}$, $C=100\,\mathrm{\mu F}$, $q(0)=2{,}0\times10^{-4}\,\mathrm C$ et $i(0)=0$. Trouver $R$ pour atteindre le régime critique et exprimer $q(t)$.


\end{document}
